\section{Minimal closed forms for second solutions}
\label{sec:minimal}

Theorem~\ref{thm:LB} needs a harmonic-monomial weight. For Ap\'ery's $b_n$ the fitting
procedure of Section~\ref{ssec:decomp} returned one of startling brevity, and this section
proves that it is correct. We then prove two more closed forms of the same kind, and
record a sharp obstruction: the method provably fails for the $\zeta(2)$-Ap\'ery
sequence.

\subsection{The minimal form}

\begin{theorem}[minimal form]\label{thm:minimal}
\Proved\ For every integer $n\ge0$,
\[
   \sum_{k=0}^{n}\bin nk^2\bin{n+k}k^2\bigl(2\HH3n-\HH3k\bigr)
   \;=\;
   \sum_{k=0}^{n}\bin nk^2\bin{n+k}k^2
     \Bigl(\HH3n+\sum_{m=1}^{k}\frac{(-1)^{m-1}}{2m^3\bin nm\bin{n+m}m}\Bigr)
   \;=\;b_n .
\]
Equivalently,
$\sum_{k=0}^{n}A(n,k)\bigl(\HH3n-\HH3k-\sum_{m=1}^ku(n,m)\bigr)=0$.
\end{theorem}

Write $B_{\min}(n)$ for the left-hand sum. The proof is by the classical route: both
sides satisfy Ap\'ery's recurrence and agree at $n=0,1$. Set
\begin{equation}\label{eq:Ldef}
   P(n):=34n^3+51n^2+27n+5,\qquad
   (Lu)(n):=(n+1)^3u(n+1)-P(n)u(n)+n^3u(n-1),
\end{equation}
and, for integers $n\ge1$, $k\ge0$, the \emph{base term}
\begin{equation}\label{eq:Phi}
   \Phi(n,k):=\frac{\bin{n+1}k^2\bin{n+k-1}k^2}{n^2(n+1)^2}.
\end{equation}

\begin{lemma}[conversion]\label{lem:conv}
\Proved\ For all integers $n\ge1$, $k\ge0$:
\begin{align}
   A(n,k)&=\Phi(n,k)(n+1-k)^2(n+k)^2, \tag{0a}\\
   A(n+1,k)&=\Phi(n,k)(n+k)^2(n+k+1)^2, \tag{0b}\\
   A(n-1,k)&=\Phi(n,k)(n+1-k)^2(n-k)^2, \tag{0c}\\
   \Phi(n,k+1)(k+1)^4&=\Phi(n,k)(n+1-k)^2(n+k)^2. \tag{0d}
\end{align}
\end{lemma}

\begin{proof}
Each reduces to the binomial conversions
$\bin nk(n+1)=\bin{n+1}k(n+1-k)$, $\bin{n+k}kn=\bin{n+k-1}k(n+k)$,
$\bin{n+k+1}kn(n+1)=\bin{n+k-1}k(n+k)(n+k+1)$,
$\bin{n-1}kn(n+1)=\bin{n+1}k(n+1-k)(n-k)$, and, for (0d),
%% FIXED [MINOR]: only the last two are literal instances of $\bin N{j+1}(j+1)=\bin Nj(N-j)$;
%% FIXED the first four move the UPPER index and are instances of the companion form. The
%% FIXED hedge present at MINIMAL_FORM_PROOF.md:48-51 had been dropped.
$\bin{n+1}{k+1}(k+1)=\bin{n+1}k(n+1-k)$, $\bin{n+k}{k+1}(k+1)=\bin{n+k-1}k(n+k)$. The last
two, those of (0d), are absorption identities $\bin N{j+1}(j+1)=\bin Nj(N-j)$ --- the second of
them followed by one step of the companion. The first four move the \emph{upper} index instead
and are instances of that companion, $\bin{N+1}j(N+1-j)=\bin Nj(N+1)$, applied once or twice.
All six hold for
\emph{every} integer $k\ge0$, including the degenerate range $k>n$ where both sides
vanish: the conversions are stated multiplicatively, so no $0/0$ arises.
\end{proof}

Lemma~\ref{lem:conv} is the device that removes every boundary subtlety: after it, all
identities below are $\Phi(n,k)$ times polynomial identities in $\Q[n,k]$.
\VerifiedR{all six conversions, $n=1..12$, $k=0..n+4$, difference exactly $0$}.

\subsection{Step B: the Zeilberger certificate for \texorpdfstring{$a_n$}{a\_n}}

Set
\begin{equation}\label{eq:TG}
   T(n,k):=4(2n+1)\bigl(2k^2-3k-4n^2-4n\bigr),
   \qquad
   G(n,k):=\Phi(n,k)\,k^4\,T(n,k).
\end{equation}

\begin{lemma}[CERT-1]\label{lem:cert1}
\Proved\ As an identity in $\Q[n,k]$,
\begin{multline*}
  T(n,k+1)(n+1-k)^2(n+k)^2-k^4T(n,k)\\
  =(n+1)^3(n+k+1)^2(n+k)^2-P(n)(n+1-k)^2(n+k)^2+n^3(n-k)^2(n+1-k)^2 .
\end{multline*}
\end{lemma}

\begin{proposition}\label{prop:B}
\Proved\ For all integers $n\ge1$, $k\ge0$,
\begin{equation}\label{eq:dagger}
   (n+1)^3A(n+1,k)-P(n)A(n,k)+n^3A(n-1,k)\;=\;G(n,k+1)-G(n,k).
\end{equation}
\end{proposition}

\begin{proof}
By (0a)--(0c) the left side is $\Phi(n,k)$ times the right-hand side of
Lemma~\ref{lem:cert1}. By (0d),
$G(n,k+1)=\Phi(n,k+1)(k+1)^4T(n,k+1)=\Phi(n,k)(n+1-k)^2(n+k)^2T(n,k+1)$ and
$G(n,k)=\Phi(n,k)k^4T(n,k)$, so the right side is $\Phi(n,k)$ times the left-hand side of
Lemma~\ref{lem:cert1}.
\end{proof}

The boundary values are $G(n,0)=0$ (factor $k^4$), $G(n,k)=0$ for $k\ge n+2$ (factor
$\bin{n+1}k^2$), and $G(n,n+1)=-4(n+1)^3\bin{2n+1}{n+1}^2$. Summing \eqref{eq:dagger} over
$k=0,\dots,N$ with $N\ge n+1$ telescopes to $G(n,N+1)-G(n,0)=0$, which recovers Ap\'ery's
recurrence $(La)(n)=0$.

\subsection{Step B\texorpdfstring{$'$}{'} and Step C: \texorpdfstring{$L$ annihilates $B_{\min}$}{L annihilates B\_min}}

Write $w(n,k):=2\HH3n-\HH3k$, so $B_{\min}(n)=\sum_{k\ge0}A(n,k)w(n,k)$ (terms with $k>n$
vanish). Since $w(n+1,k)=w(n,k)+2/(n+1)^3$ and $w(n-1,k)=w(n,k)-2/n^3$,
\[
   (LB_{\min})(n)=\sum_{k\ge0}w(n,k)\bigl(G(n,k+1)-G(n,k)\bigr)+2a_{n+1}-2a_{n-1},
\]
by Proposition~\ref{prop:B}. Abel summation, with $w(n,k+1)-w(n,k)=-1/(k+1)^3$ and the
boundary values of $G$, turns this into
\begin{equation}\label{eq:Bprime}
   (LB_{\min})(n)\;=\;\sum_{j\ge1}\frac{G(n,j)}{j^3}+2a_{n+1}-2a_{n-1}.
\end{equation}

Now set
\begin{equation}\label{eq:UK}
   U(n,k):=-4(k-1)(2n+1)\bigl(2n^2+2n+1-k\bigr),
   \qquad
   K(n,k):=\frac{\Phi(n,k)k^4U(n,k)}{n^2(n+1)^2}.
\end{equation}

\begin{lemma}[CERT-2]\label{lem:cert2}
\Proved\ As an identity in $\Q[n,k]$,
\begin{multline*}
   U(n,k+1)(n+1-k)^2(n+k)^2-k^4U(n,k)\\
   =n^2(n+1)^2\Bigl[kT(n,k)+2\bigl((n+k)^2(n+k+1)^2-(n+1-k)^2(n-k)^2\bigr)\Bigr],
\end{multline*}
and the bracket equals $4k(2n+1)(4k^2-3k-2n-2n^2)$.
\end{lemma}

\begin{proposition}\label{prop:C}
\Proved\ For all integers $n\ge1$, $k\ge0$, with the convention $G(n,0)/0^3:=0$
(consistent, since $G(n,k)/k^3=\Phi(n,k)kT(n,k)$),
\begin{equation}\label{eq:ddagger}
   \frac{G(n,k)}{k^3}+2A(n+1,k)-2A(n-1,k)\;=\;K(n,k+1)-K(n,k).
\end{equation}
Consequently $\sum_{j\ge1}G(n,j)/j^3+2a_{n+1}-2a_{n-1}=0$, since $K(n,0)=0$ and
$K(n,k)=0$ for $k\ge n+2$.
\end{proposition}

\begin{proof}
By (0b), (0c) the left side is $\Phi(n,k)\bigl[kT(n,k)
+2((n+k)^2(n+k+1)^2-(n+1-k)^2(n-k)^2)\bigr]$; by (0d) the right side is
$\Phi(n,k)\bigl[(n+1-k)^2(n+k)^2U(n,k+1)-k^4U(n,k)\bigr]/(n^2(n+1)^2)$. Equality is
Lemma~\ref{lem:cert2}. Summing over $k=0,\dots,N$ with $N\ge n+1$ telescopes.
\end{proof}

\begin{theorem}\label{thm:LBmin}
\Proved\ $(LB_{\min})(n)=0$ for all $n\ge1$.
\end{theorem}

\begin{proof}
Combine \eqref{eq:Bprime} with Proposition~\ref{prop:C}.
\end{proof}

\subsection{Step D: the classical companion satisfies the same recurrence}

This is the half in which Ap\'ery's harmonic weight $c(n,k)$ must be differenced
\emph{in $n$}. Everything collapses because of two one-line Gosper certificates whose
entire content is the pair of polynomial identities
\begin{equation}\label{eq:miracle}
   m^2+(n+m)(n-m)=n^2,
   \qquad
   m^2+(n+1-m)(n+1+m)=(n+1)^2 .
\end{equation}

Put $D(n,k):=\bin nk\bin{n+k}k$ (so $A=D^2$), and
\[
 \begin{gathered}
   \varphi(n,m):=m\,u(n,m)=\frac{(-1)^{m-1}((m-1)!)^2(n-m)!}{2(n+m)!},\\
   \Lambda(n,k):=(n-k)\varphi(n,k+1)=\frac{(-1)^k}{2(n+k+1)D(n,k)} .
 \end{gathered}
\]
Since $D(n,k)\ne0$ for $0\le k\le n$, every division below is legitimate on its stated
range.

\begin{lemma}[the two inner Gosper certificates]\label{lem:D1}
\Proved\ For $1\le m\le n$ put
\[
   Y_1(m):=-\frac{(n+m)\varphi(n,m)}{n^2},
   \qquad
   Y_2(m):=-\frac{(n+1-m)\varphi(n,m)}{(n+1)^2} .
\]
Then, using $\varphi(n,m+1)/\varphi(n,m)=-m^2/\bigl((n-m)(n+m+1)\bigr)$,
\[
   Y_1(m+1)-Y_1(m)=\frac{\varphi(n,m)}{n-m}\ (1\le m\le n-1),
   \quad
   Y_2(m+1)-Y_2(m)=\frac{\varphi(n,m)}{n+m+1}\ (1\le m\le n),
\]
the two brackets collapsing by the identities \eqref{eq:miracle}.
\end{lemma}

\begin{lemma}[the $n$-differences of Ap\'ery's weight]\label{lem:D2}
\Proved\ Using $u(n+1,m)=u(n,m)\frac{n+1-m}{n+m+1}$ and $u(n-1,m)=u(n,m)\frac{n+m}{n-m}$,
i.e.\ $u(n+1,m)-u(n,m)=-2\varphi(n,m)/(n+m+1)$ and
$u(n-1,m)-u(n,m)=2\varphi(n,m)/(n-m)$, Lemma~\ref{lem:D1} telescopes with
$Y_2(1)=-\tfrac1{2(n+1)^3}$, $Y_1(1)=-\tfrac1{2n^3}$ to give
\[
 \begin{gathered}
   c(n+1,k)-c(n,k)=\frac{2\Lambda(n,k)}{(n+1)^2}\quad(0\le k\le n),\\
   c(n-1,k)-c(n,k)=-\frac{2(n+k+1)\Lambda(n,k)}{(n-k)n^2}\quad(0\le k\le n-1).
 \end{gathered}
\]
\end{lemma}

\begin{remark}[Ap\'ery's miracle]\label{rem:miracle}
In Lemma~\ref{lem:D2} the rational pieces $1/(n+1)^3$ and $-1/n^3$ coming from
$\HH3{n\pm1}-\HH3n$ cancel \emph{exactly} against the boundary values $Y_2(1)$, $Y_1(1)$
of the inner Gosper antidifferences, and that cancellation is precisely the polynomial
identity $m^2+(n+1-m)(n+m+1)=(n+1)^2$. This is why Ap\'ery's weight carries the $\tfrac12$,
the $(-1)^{m-1}$ and the $m^3\bin nm\bin{n+m}m$: the weight is not found, it is
\emph{forced} --- $c(n,k)$ is the unique weight, up to the kernel of the recurrence, whose
$n$-difference is a \emph{single hypergeometric term} rather than a sum.
\end{remark}

\begin{lemma}[the correction term]\label{lem:D3}
\Proved\ For $0\le k\le n$ set
$\mathcal E(n,k):=(n+1)^3A(n+1,k)\bigl(c(n+1,k)-c(n,k)\bigr)+n^3A(n-1,k)\bigl(c(n-1,k)-c(n,k)\bigr)$,
the second summand being $0$ at $k=n$ because $A(n-1,n)=0$. Then by Lemma~\ref{lem:D2},
(0b), (0c) and $A(n,k)\Lambda(n,k)=(-1)^kD(n,k)/\bigl(2(n+k+1)\bigr)$,
\[
   \mathcal E(n,k)=(-1)^kD(n,k)\Bigl[\frac{(n+1)(n+k+1)}{(n+1-k)^2}-\frac{n(n-k)}{(n+k)^2}\Bigr].
\]
\end{lemma}

\begin{proposition}\label{prop:D4}
\Proved\ For $0\le k\le n$, with $B(n-1,k):=0$ for $k\ge n$ and $B(n,k):=A(n,k)c(n,k)$,
\[
   (n+1)^3B(n+1,k)-P(n)B(n,k)+n^3B(n-1,k)
   =c(n,k)\bigl(G(n,k+1)-G(n,k)\bigr)+\mathcal E(n,k).
\]
At $k=n+1$ only the first term survives: the left side equals
$(n+1)^3A(n+1,n+1)c(n+1,n+1)$.
\end{proposition}

\begin{proof}
Substitute $c(n\pm1,k)=c(n,k)+\bigl(c(n\pm1,k)-c(n,k)\bigr)$ and use
Proposition~\ref{prop:B}.
\end{proof}

\begin{lemma}[the top term]\label{lem:D5}
\Proved\ $G(n,n+1)=-4(n+1)^3\bin{2n+1}{n+1}^2=-(n+1)^3A(n+1,n+1)$, using
$\bin{2n+2}{n+1}=2\bin{2n+1}{n+1}$. Moreover
$u(n+1,n+1)=\frac{(-1)^n}{4(n+1)^3\bin{2n+1}{n+1}}$ and
$\Lambda(n,n)=\frac{(-1)^n}{2(n+1)\bin{2n+1}{n+1}}$, so that
$c(n+1,n+1)-c(n,n)=\tfrac54\,(-1)^n\bigl/\bigl((n+1)^3\bin{2n+1}{n+1}\bigr)$ and
\[
   \mathrm{Top}(n):=c(n,n)G(n,n+1)+(n+1)^3A(n+1,n+1)c(n+1,n+1)
   =5(-1)^n\bin{2n+1}{n+1}.
\]
\end{lemma}

This is where the classical $5$ of $\zeta(3)=\tfrac52\sum_m(-1)^{m-1}m^{-3}\bin{2m}m^{-1}$
enters.

Summing Proposition~\ref{prop:D4} over $k=0,\dots,n$, adding the $k=n+1$ term, and
Abel-summing $\sum_{k=0}^nc(n,k)(G(n,k+1)-G(n,k))=c(n,n)G(n,n+1)-\sum_{j=1}^nu(n,j)G(n,j)$
(using $G(n,0)=0$), we get
\begin{equation}\label{eq:Lb-assembly}
   (Lb)(n)=\mathrm{Top}(n)+\sum_{k=0}^{n}V(n,k),
   \qquad V(n,k):=\mathcal E(n,k)-u(n,k)G(n,k),
\end{equation}
the $k=0$ term of the second piece being $0$. Since
\[
  u(n,k)G(n,k)=\frac{(-1)^{k-1}D(n,k)\,k\,T(n,k)}{2(n+1-k)^2(n+k)^2},
\]
we have $V(n,k)=(-1)^kD(n,k)W(n,k)$ with
\[
   W(n,k)=\frac{(n+1)(n+k+1)}{(n+1-k)^2}-\frac{n(n-k)}{(n+k)^2}
          +\frac{kT(n,k)}{2(n+1-k)^2(n+k)^2}.
\]

\begin{lemma}[CERT-3]\label{lem:cert3}
\Proved\ As an identity in $\Q[n,k]$,
\[
   2(n+1)(n+k+1)(n+k)^2-2n(n-k)(n+1-k)^2+kT(n,k)+10k(2n+1)(n+1-k)(n+k)=0,
\]
that is,
\[
 \begin{gathered}
   W(n,k)=\frac{-5k(2n+1)}{(n+1-k)(n+k)},\qquad\text{whence}\\
   V(n,k)=\frac{-5(2n+1)}{n(n+1)}\,(-1)^k\,k\,\bin{n+1}k\bin{n+k-1}k,
 \end{gathered}
\]
using $\bin nk/(n+1-k)=\bin{n+1}k/(n+1)$ and $\bin{n+k}k/(n+k)=\bin{n+k-1}k/n$.
\end{lemma}

The entire weight-three structure thus collapses to a single linear-over-quadratic
rational function.

\begin{lemma}[residual binomial identity]\label{lem:D6}
\Proved\
\[
   \text{(D-BIN)}\qquad
   \sum_{k=0}^{n}(-1)^k\,k\,\bin{n+1}k\bin{n+k-1}k\;=\;(-1)^n\,n\bin{2n}n .
\]
\end{lemma}

\begin{proof}
Put $s(n,k):=(-1)^kk\bin{n+1}k\bin{n+k-1}k$ and $Y(n,k):=-k(k-1)s(n,k)/(n(n+1))$. The
absorption identities $\bin{n+1}{k+1}(k+1)=\bin{n+1}k(n+1-k)$ and
$\bin{n+k}{k+1}(k+1)=\bin{n+k-1}k(n+k)$ of Lemma~\ref{lem:conv} give
$(k+1)s(n,k+1)=-(-1)^k\bin{n+1}k\bin{n+k-1}k(n+1-k)(n+k)$, hence for \emph{every} integer
$k\ge0$, division-free,
\begin{multline*}
   n(n+1)\bigl(Y(n,k+1)-Y(n,k)\bigr)\\
   =(-1)^kk\bin{n+1}k\bin{n+k-1}k\bigl[(n+1-k)(n+k)+k(k-1)\bigr]=n(n+1)s(n,k),
\end{multline*}
by the polynomial identity $(n+1-k)(n+k)+k(k-1)=n(n+1)$. So $Y(n,\cdot)$ telescopes
$s(n,\cdot)$ and $\sum_{k=0}^ns(n,k)=Y(n,n+1)-Y(n,0)=Y(n,n+1)=-s(n,n+1)
=(-1)^nn\bin{2n}n$, using $\bin{2n}{n+1}=n\bin{2n}n/(n+1)$.
\end{proof}

\begin{theorem}\label{thm:Lb}
\Proved\ $(Lb)(n)=0$ for all $n\ge1$.
\end{theorem}

\begin{proof}
By \eqref{eq:Lb-assembly}, Lemma~\ref{lem:cert3} and Lemma~\ref{lem:D6},
\begin{align*}
   \sum_{k=0}^{n}V(n,k)
   &=\frac{-5(2n+1)}{n(n+1)}(-1)^nn\bin{2n}n
    =-5(-1)^n\bin{2n}n\frac{2n+1}{n+1}\\
   &=-5(-1)^n\bin{2n+1}{n+1}=-\mathrm{Top}(n). \qedhere
\end{align*}
\end{proof}

\begin{proof}[Proof of Theorem~\ref{thm:minimal}]
Theorems~\ref{thm:LBmin} and \ref{thm:Lb} say that both sides are annihilated by $L$,
whose leading coefficient $(n+1)^3$ never vanishes for $n\ge1$. It remains to check
agreement at $n=0,1$. We have $B_{\min}(0)=A(0,0)(2\HH30-\HH30)=0=b_0$, and
$B_{\min}(1)=A(1,0)\cdot2+A(1,1)\cdot1=1\cdot2+4\cdot1=6$; on the other side
$u(1,1)=1/(2\cdot1^3\bin11\bin21)=1/4$, so $c(1,0)=1$, $c(1,1)=5/4$ and
$b_1=1\cdot1+4\cdot\tfrac54=6$. Induct.
\end{proof}

\begin{table}[ht]
\centering\small
\caption{Certificate inventory for Theorem~\ref{thm:minimal}. All are verified by exact
polynomial arithmetic (\texttt{Expand}$\to0$).}
\label{tab:certs}
\begin{tabular}{@{}llp{6.4cm}@{}}
\toprule
\# & role & certificate\\
\midrule
CERT-1 & Zeilberger certificate for $a_n$ & $T(n,k)=4(2n+1)(2k^2-3k-4n^2-4n)$, $G=\Phi k^4T$\\
CERT-2 & Gosper certificate for the $w$-correction & $U(n,k)=-4(k-1)(2n+1)(2n^2+2n+1-k)$, $K=\Phi k^4U/(n^2(n+1)^2)$\\
CERT-3 & the $W$-collapse & $W(n,k)=-5k(2n+1)/\bigl((n+1-k)(n+k)\bigr)$\\
CERT-4 & inner Gosper certificate, $\Delta^+_n$ & $Y_2(m)=-(n+1-m)\varphi(n,m)/(n+1)^2$\\
CERT-5 & inner Gosper certificate, $\Delta^-_n$ & $Y_1(m)=-(n+m)\varphi(n,m)/n^2$\\
CERT-6 & Gosper certificate for (D-BIN) & $Y(n,k)=-k(k-1)s(n,k)/\bigl(n(n+1)\bigr)$\\
\bottomrule
\end{tabular}
\end{table}

The supporting polynomial identities, all \texttt{Expand}$\to0$, are
$m^2+(n+m)(n-m)=n^2$; $m^2+(n+1-m)(n+m+1)=(n+1)^2$; $(n+1-k)(n+k)+k(k-1)=n(n+1)$; and
CERT-1, CERT-2, CERT-3 as displayed. Numerically, $B_{\min}(n)-b_n=0$ exactly for
$n=0..50$ in exact rational arithmetic, in two independent implementations
\VerifiedR{$n\le50$}; and $B_{\min}$ reproduces the literature values
$0,6,\tfrac{351}4,\tfrac{62531}{36},\tfrac{11424695}{288},\tfrac{35441662103}{36000}$.

\begin{remark}[formalisation-friendliness]\label{rem:leanfriendly}
Nothing above uses Gamma functions, limits or $0/0$: Lemma~\ref{lem:conv} converts every
binomial ratio into a \emph{multiplicative} absorption identity valid for all integers
$k\ge0$, after which each step is a polynomial identity in $\Q[n,k]$ plus a finite
telescoping sum. The only analytic input is $\HH3{n+1}-\HH3n=1/(n+1)^3$.
\end{remark}

\subsection{Two more closed forms, and the mechanism behind them}

\begin{observation}[symmetrisation]\label{obs:symm}
\Proved\ For the $d$-fold binomial sums $A^{(d)}(n)=\sum_k\bin nk^d$ the summand is
symmetric under $k\mapsto n-k$, so the fitted weights of Table~\ref{tab:decomp} may be
replaced by their $k\leftrightarrow n-k$ averages without changing the sum:
\[
   \tfrac1{d+1}\HH2k+\tfrac d{d+1}H_k(H_k-H_{n-k})
   \;\rightsquigarrow\;
   \tfrac1{2(d+1)}\bigl(\HH2k+\HH2{n-k}\bigr)+\tfrac d{2(d+1)}\bigl(H_k-H_{n-k}\bigr)^2 .
\]
\VerifiedR{exactly, $n=0..12$, $d=3,4$}. Concretely,
\[
 \begin{gathered}
   B_{\mathbf A}(n)=\sum_k\bin nk^3\Bigl[\tfrac18\bigl(\HH2k+\HH2{n-k}\bigr)+\tfrac38\bigl(H_k-H_{n-k}\bigr)^2\Bigr],\\
   B_{s_{10}}(n)=\sum_k\bin nk^4\Bigl[\tfrac1{10}\bigl(\HH2k+\HH2{n-k}\bigr)+\tfrac25\bigl(H_k-H_{n-k}\bigr)^2\Bigr].
 \end{gathered}
\]
\end{observation}

This is not cosmetic. In the symmetric form the only weight-one letter that ever appears
is the single antisymmetric combination $\delta(n,k):=H_k-H_{n-k}$, because
$\sigma:=\HH2k+\HH2{n-k}$ differences to \emph{rational} functions under both $n\mapsto
n\pm1$ and $k\mapsto k+1$. The certificate module is therefore of rank $2$
($\{1,\delta\}$) instead of rank $4$ ($\{1,H_k,H_{n-k},H_n\}$), and one faces two Gosper
problems instead of four.

The template is then: (i) a Zeilberger certificate
$G(n,k)=S(n,k)k^d\rho(n,k)/(n+1-k)^d$, whose defining equation collapses to one polynomial
equation in $\Q[n,k]$ because $S(n,k+1)/S(n,k)=((n-k)/(k+1))^d$, so that the shifted
certificate is the summand times a polynomial; (ii) a split of the recurrence applied to
$S\cdot w$ into $w\cdot\Delta_kG$ plus a correction whose $n$-differences lie in
$\Q(n,k)+\Q(n,k)\delta$; (iii) Abel summation; (iv) two Gosper problems
\[
 \begin{gathered}
   r(n,k)z_1(n,k+1)-z_1(n,k)=B(n,k),\\
   r(n,k)z_0(n,k+1)-z_0(n,k)=A(n,k)-r(n,k)z_1(n,k+1)e(n,k),
 \end{gathered}
\]
with $r=S(n,k+1)/S(n,k)$ and $e=\frac1{k+1}+\frac1{n-k}$, both solved by rational ans\"atze
$z_1=N(k)/(n+1-k)^{p}$, $z_0=M(k)/(n+1-k)^{p+1}$; and (v) three scalar boundary
identities, of which the first,
\[
   \text{(b1)}\qquad \rho(n,n+1)=-c_+(n)\,S(n+1,n+1)/S(n,n),
\]
is structural: it says that the Zeilberger certificate at the top of its support
reproduces the leading term of the recurrence, and it is what forces a harmonic
decomposition to exist at all.

%% FIXED [MAJOR-2]: \Proved\ is defined at sec-intro.tex as "a complete proof is written out
%% FIXED here", and neither this theorem nor Theorem \ref{thm:s10} writes one out: the Gosper
%% FIXED equations the certificates solve are not stated, and Theorem \ref{thm:s10}'s N_10, M_10
%% FIXED are not in the paper at all. Both are now labelled for what they are: certificates
%% FIXED exhibited and checked, with the verification exact.
\begin{theorem}[Franel]\label{thm:franel}
\Proved\ modulo the certificate check described after Theorem~\ref{thm:s10}; the displayed
identities are \Proved\ (\texttt{Factor}$\to0$, exact).
With $A(n)=\sum_k\bin nk^3$ (A000172), $(n+1)^2u_{n+1}=(7n^2+7n+2)u_n+8n^2u_{n-1}$
and $B(n)=\sum_k\bin nk^3\bigl[\tfrac18(\HH2k+\HH2{n-k})+\tfrac38(H_k-H_{n-k})^2\bigr]$,
one has $(LB)=0$ for the Franel operator, so $B$ is the second solution, with values
$B(0),\dots,B(8)$ equal to
$0$, $1$, $4$, $\tfrac{208}9$, $\tfrac{1280}9$, $\tfrac{208384}{225}$,
$\tfrac{1404928}{225}$, $\tfrac{95174656}{2205}$, $\tfrac{3351248896}{11025}$.
The certificates are
\[
 \begin{gathered}
  \rho_F(n,k)=-\tfrac1n\bigl(4-12k+12k^2-4k^3+22n-39kn+18k^2n+32n^2-27kn^2+14n^3\bigr),\\
  G_F=\tfrac{k^3\rho_F\bin{n+1}k^3}{(n+1)^3},
 \end{gathered}
\]
$z_1=N_F(k)/(n+1-k)^4$, $z_0=M_F(k)/(n+1-k)^5$ with
{\small
\begin{align*}
  N_F&=\tfrac{3k^2}{4n}\bigl(4-16k+24k^2-16k^3+4k^4+26n-68kn+63k^2n-20k^3n\\
     &\qquad\qquad+54n^2-88kn^2+39k^2n^2+46n^3-36kn^3+14n^4\bigr),\\
  M_F&=-\tfrac{k}{2n}\bigl(2-10k+20k^2-20k^3+10k^4-2k^5+15n-50kn+70k^2n-45k^3n+11k^4n\\
     &\qquad\qquad+40n^2-90kn^2+80k^2n^2-25k^3n^2+50n^3-70kn^3+30k^2n^3\\
     &\qquad\qquad+30n^4-20kn^4+7n^5\bigr),
\end{align*}}
and the boundary identities, all \texttt{Factor}$\to0$, are
$\rho_F(n,n+1)=-(n+1)^2$, $N_F(n)=\tfrac34\bigl((n+1)^5-(n+1)\bigr)$,
$M_F(n)=-\tfrac12\bigl(1+(n+1)^5\bigr)$.
\end{theorem}

\begin{theorem}[$s_{10}=\sum_k\bin nk^4$]\label{thm:s10}
\Proved\ modulo the certificate check described below.
With $(n+1)^3u_{n+1}=(2n+1)(6n^2+6n+2)u_n+n(64n^2-4)u_{n-1}$ and
$B(n)=\sum_k\bin nk^4\bigl[\tfrac1{10}(\HH2k+\HH2{n-k})+\tfrac25(H_k-H_{n-k})^2\bigr]$,
one has $(LB)=0$, with $B(0),\dots,B(6)$ equal to
$0$, $1$, $\tfrac{21}4$, $\tfrac{1001}{18}$, $\tfrac{85085}{144}$,
$\tfrac{4203199}{600}$, $\tfrac{52055003}{600}$.
The certificate is $G_{10}=k^4\rho_{10}\bin{n+1}k^4/(n+1)^4$ with
{\small
\begin{align*}
 \rho_{10}(n,k)=-\tfrac1{n^3}\bigl(&-4k+16k^2-24k^3+16k^4-4k^5\\
 &+10n-72kn+172k^2n-184k^3n+90k^4n-16k^5n\\
 &+80n^2-368kn^2+600k^2n^2-416k^3n^2+104k^4n^2\\
 &+255n^3-796kn^3+818k^2n^3-276k^3n^3\\
 &+385n^4-756kn^4+374k^2n^4+275n^5-260kn^5+75n^6\bigr),
\end{align*}}
$z_1=N_{10}(k)/(n+1-k)^5$ and $z_0=M_{10}(k)/(n+1-k)^6$ (both of degree $9$ in $k$, with
factors $k^3$ and $k^2$; full coefficients in the certificate file, see
Section~\ref{sec:data}), and boundary identities $\rho_{10}(n,n+1)=-(n+1)^3$,
$N_{10}(n)=\tfrac45\bigl((n+1)^7-(n+1)^2\bigr)$,
$M_{10}(n)=-\tfrac12\bigl((n+1)+(n+1)^7\bigr)$, all \texttt{Factor}$\to0$.
\end{theorem}

\begin{remark}[what the two theorems above do and do not contain]\label{rem:certstatus}
Both are proved by the same route as Theorem~\ref{thm:Lb}: one exhibits $\rho$, $G$,
$z_1$, $z_0$, checks $\Delta_k G=\rho S$ and the two Gosper identities as polynomial identities
in $\Q(n)[k]$, and checks the boundary identities. Every displayed identity above was verified
that way, by exact expansion to $0$. What is \emph{not} written out here is the derivation: we
do not state the Gosper key equations that $z_1,z_0$ solve, and for Theorem~\ref{thm:s10} the
polynomials $N_{10}$, $M_{10}$ themselves are not printed --- their full coefficients are in the
certificate file (Section~\ref{sec:data}). A reader wanting a self-contained proof must read
that file. We flag this because \Proved\ is defined in Section~\ref{ssec:evidence} as ``a
complete proof is written out here'', and for these two theorems the certificate half is
delegated rather than printed.
\end{remark}

\begin{observation}[the pattern]\label{obs:pattern}
With $\beta_d=\frac d{2(d+1)}$ the $\delta^2$-coefficient and $c_+(n)=(n+1)^{d-1}$, both
theorems read
\[
 \begin{gathered}
   \rho(n,n+1)=-(n+1)^{d-1},\qquad
   N(n)=2\beta_d\bigl((n+1)^{2d-1}-(n+1)^{d-2}\bigr),\\
   M(n)=-\tfrac12\bigl((n+1)^{2d-1}+(n+1)^{d-3}\bigr).
 \end{gathered}
\]
It is natural to conjecture that this is the general-$d$ shape, i.e.\ that
Chamberland--Straub's conjecture that $A^{(d)}$ has Ap\'ery limit $\zeta(2)/(d+1)$ is
provable this way for every $d$ for which the sequence satisfies an order-two recurrence
\VerifiedR{$d=3,4$ only}.
\end{observation}

\subsection{A sharp obstruction: the \texorpdfstring{$\zeta(2)$}{zeta(2)}-Ap\'ery case}
\label{ssec:obstruction}

The template of the previous subsection does not merely fail to apply to the
$\zeta(2)$-Ap\'ery sequence $\mathbf D$; it provably cannot apply.

Let $A(n)=\sum_k\bin nk^2\bin{n+k}k$ (A005258: $1,3,19,147,1251,\dots$),
$(n+1)^2u_{n+1}=(11n^2+11n+3)u_n+n^2u_{n-1}$, and
$B(n)=\sum_k\bin nk^2\bin{n+k}k\cdot\frac15\bigl[\HH2n+H_k(2H_k-H_{n-k}-H_n)\bigr]$, with
$(LB)=0$ \VerifiedR{$n\le25$}. Everything up to the Gosper step goes through, and
beautifully: the Zeilberger certificate is the smallest in this paper,
\[
   \rho_{\mathbf D}(n,k)=k^2+k(1+6n)-4-15n-11n^2,
   \qquad
   G_{\mathbf D}(n,k)=\frac{k^3\rho_{\mathbf D}(n,k)\bin{n+1}k^2\bin{n+k-1}k}{n(n+1)^2},
\]
\Proved\ (\texttt{Expand}$\to0$; \VerifiedR{pointwise, $n=1..9$, $k=0..n+3$}), and the
structural boundary identity (b1) holds exactly:
\[
   \rho_{\mathbf D}(n,n+1)=-2(n+1)(2n+1)=-(n+1)^2\frac{\bin{2n+2}{n+1}}{\bin{2n}n}
   =-c_+(n)\frac{S(n+1,n+1)}{S(n,n)} .
\]
The correction decomposes as $Y(n,k)/S(n,k)=A+B_1H_k+B_2(H_{n-k}+H_n)$ with
$B_2=\rho_{\mathbf D}(n,k+1)/\bigl(5(k+1)\bigr)$ \VerifiedR{exactly, $n=2..8$,
$0\le k\le n-1$}.

\begin{proposition}[the $H_{n-k}$ component is not Gosper-summable]\label{prop:obstruction}
\Proved\ Let $T(k):=S_{\mathbf D}(n,k)B_2(n,k)$. In Gosper's decomposition
$T(k+1)/T(k)=\frac{a(k)}{b(k)}\frac{c(k+1)}{c(k)}$ one has $a=(n-k)^2(n+k+1)$,
$b=(k+1)^2(k+2)$, $c(k)=\rho_{\mathbf D}(n,k+1)$, and $\gcd\bigl(a(k),b(k+j)\bigr)=1$ for
all $j\ge0$ (the roots are $n,-n-1$ versus $-j-1,-j-2$, with $n$ generic). Gosper's key
equation is $a(k)x(k+1)-b(k-1)x(k)=c(k)$. Since
\[
   a(k)-b(k-1)=-n\bigl(k^2+(n+2)k-n(n+1)\bigr)
\]
has degree $2$ with leading coefficient $-n\ne0$, the degree bound forces $\deg x=0$; and
$x=x_0$ constant forces $x_0=-1/n$ from the $k^2$ coefficient, which then produces the
$k$-coefficient $n+2$ against the required $3+6n$. Contradiction. Independently
\Verified: the corresponding linear system has no solution for any $\deg x\le8$.
\end{proposition}

\begin{remark}[diagnosis]\label{rem:diagnosis}
In the Franel and $s_{10}$ cases the weight is $k\leftrightarrow n-k$ symmetric, the
certificate module has rank $2$, and each component telescopes. Here the weight is
genuinely three-lettered ($H_k$, $H_{n-k}$, $H_n$), the components do not telescope
separately, and the weight-two part of any enlarged ansatz is forced to vanish
($\Delta(Sz)=0\Rightarrow z=0$). Closing this case needs either a different gauge for $w$
--- the harmonic-monomial fit has a large kernel, $11$-dimensional in the $\zeta(3)$ case,
and a kernel shift changes the antidifference --- or a genuine $\Pi\Sigma$-field ansatz
allowing nested sums rather than rational $\times$ harmonic-monomial. This is \Open, and
it is the first place in this programme where the rank-two mechanism is provably
insufficient.
\end{remark}

\subsection{Literature}
\label{ssec:minlit}

The recurrence $(Lb)=0$ itself is not new. It is proved in full by Schneider
\cite{Schneider2007} using Karr's $\Pi\Sigma$ difference fields and creative
telescoping, with a certificate whose numerator, after the substitution $n\mapsto n-1$,
is exactly the numerator of our $G(n,k)$ --- an independent confirmation of the
normalisation of \eqref{eq:TG}. Historically it goes back to the Zagier--Cohen
reorganisation of Ap\'ery's proof reported by van der Poorten
\cite[\S8]{vdPoorten1979}, who gives
$B_{n,k}=4(2n+1)\bigl(k(2k+1)-(2n+1)^2\bigr)\bin nk^2\bin{n+k}k^2$ and, for the second
solution, $A_{n,k}=B_{n,k}c_{n,k}+\bigl[5(2n+1)(-1)^{k-1}k/(n(n+1))\bigr]\bin nk\bin{n+k}k$
--- precisely our Step-D certificate, with $G(n,k)=B_{n,k-1}$ and $V$ as in
Lemma~\ref{lem:cert3}. But van der Poorten dispatches the verification as ``after some
massive reorganisation (9) becomes $A_{n,k}-A_{n,k-1}$'': it is exhibited, not written
out; Fischler \cite[\S1.2]{Fischler2004} likewise calls it ``une simple
v\'erification''. Zudilin \cite[eqs.~(5)--(6)]{Zudilin0202159} has the rational-function
form $s_n(t)=4(2n+1)(-2t^2+t+(2n+1)^2)$, the same polynomial under $t\leftrightarrow-k$.

What appears to be new is the \emph{minimal form} itself. We could not locate
$2\HH3n-\HH3k$ anywhere: not in OEIS \texttt{A059415} (which is exactly this $b_n$, and
whose only recorded closed form is the classical double sum), nor in
\cite{ChamStraub,Gorodetsky,Straub1401,Cooper2302}, nor in
\cite{vdPoorten1979,Schneider2007,Fischler2004,
PauleSchneider2003}. The nearest published relative is Nesterenko's Lemma~1
\cite{Nesterenko1996}, reproduced as formula (3) of \cite{Zudilin0202159},
which amounts to the weight $\HH3k+(2H_k-H_{n+k}-H_{n-k})\HH2k$ --- a longer
representation, not the minimal one. (The harmonic-monomial fitting system for $b_n$ has
an $11$-dimensional kernel, so such representations are far from unique; the pivot choice
delivered the shortest.) Similarly, (D-BIN) of Lemma~\ref{lem:D6} we could not locate as a
named identity, though it is one line from Chu--Vandermonde: from
$\sum_{k\ge0}(-1)^k\bin{n+1}k\bin{n+k-1}kx^k={}_2F_1(-(n+1),n;1;x)$, apply $x\,d/dx$ at
$x=1$ to get $-n(n+1)\,{}_2F_1(-n,n+1;2;1)=0$ for $n\ge1$, and peel off the $k=n+1$ term.
Its relative $\sum_k(-1)^k\bin nk\bin{n+k}k=(-1)^n$ is $P_n(-1)$. Finally, Paule and
Schneider \cite{PauleSchneider2003} prove Ahlgren's
$\sum_j(1-4jH_j+4jH_{n-j})\bin nj^4=(-1)^n\bin{2n}n$ --- the same right-hand side and the
same $d=4$ summand as Theorem~\ref{thm:s10} --- but only weight-one identities;
$H^{(3)}$ never appears there, so Theorem~\ref{thm:s10} is not covered by that work.
